\section{符号说明}

本文所用主要符号如表~\ref{tab:notation} 所示，其余符号将在首次出现时说明。

\begin{table}[H]
    \centering
    \caption{主要符号说明}
    \label{tab:notation}
    \begin{tabular}{@{}cl@{}}
        \toprule
        符号 & 含义 \\
        \midrule
        $\boldsymbol{z}_k^{(i)}$ & 第 $i$ 路在 $k$ 时刻的位置观测 $(x,y)^\top$ \\
        $\boldsymbol{p}(t)$ & 真实轨迹位置 \\
        $\Delta\tau$ & 方式2相对方式1的时间偏差（s）\\
        $\Delta x, \Delta y$ & 方式2的空间系统偏差（m）\\
        $S_i(t)$ & 第 $i$ 路的三次样条插值 \\
        $\boldsymbol{x}_k$ & 卡尔曼状态向量 $(p_x, p_y, v_x, v_y, a_x, a_y)^\top$ \\
        $\boldsymbol{F}, \boldsymbol{Q}, \boldsymbol{H}, \boldsymbol{R}$ & 状态转移、过程噪声、观测、观测噪声矩阵 \\
        $\sigma_i$ & 第 $i$ 路观测噪声标准差（m）\\
        $q_a$ & 过程噪声加速度抖动强度（m/s$^2$）\\
        $d$ & 机器人与目标间距离（m）\\
        $v, a$ & 机器人线速度（m/s）、加速度（m/s$^2$）\\
        $\Delta\theta$ & 同一拍照目标两次执行的方向角差（$^\circ$）\\
        $\Delta_{\text{prep}}$ & 准备时长（射击 1.5\,s，拍照 0.5\,s）\\
        \bottomrule
    \end{tabular}
\end{table}
